\documentclass[11pt,a4paper]{article}
\usepackage{amsmath,amssymb,amsthm,amsfonts}
\usepackage[margin=1in]{geometry}
\usepackage{booktabs}
\usepackage{hyperref}

\newtheorem{prop}{Proposition}
\newtheorem{lem}[prop]{Lemma}
\newtheorem{coro}[prop]{Corollary}
\newtheorem{rem}[prop]{Remark}

\def\R{\mathbb{R}}
\def\N{\mathbb{N}}
\def\E{\mathbb{E}}
\def\P{\mathbb{P}}
\def\T{\mathcal{T}}
\def\d{\,\mathrm{d}}
\def\xib{\boldsymbol{\xi}}
\def\Psib{\boldsymbol{\Psi}}
\def\Phib{\boldsymbol{\Phi}}
\def\Yb{\boldsymbol{Y}}
\def\Zb{\mathbf{Z}}
\def\Mb{\mathbf{M}}

\title{Exact simulation of L\'evy random-bridges without the martingale $\{\Mb_t\}$\\
\large Gamma and Poisson cases, with a proposed head-to-head experiment}
\author{Working note for internal review}
\date{\today}

\begin{document}
\maketitle

\begin{abstract}
Section 4 of the manuscript derives the non-anticipative representation
$\xib^{\boldsymbol{x}}_t \overset{\text{law}}{=} \int_0^t \frac{\E[\Yb\mid\xib^{\boldsymbol{x}}_s]-\xib^{\boldsymbol{x}}_s}{T-s}\d s + \Mb_t$
and states that a practical algorithm awaits an explicit form for the martingale
$\{\Mb_t\}$. This note observes that the simulation does not require $\{\Mb_t\}$ at
all. Sampling the \emph{bridge transition kernel} $\nu_{s,t,T}$ --- already derived in
the manuscript --- advances the state exactly, with no time-discretisation error and
no terminal singularity. For the gamma process the kernel is a scaled Beta law; for
the Poisson process it is a Binomial thinning. Both are elementary to simulate.
This converts Section 4 from a theoretical outlook into a working algorithm, and
unifies the Gaussian, gamma and Poisson cases under one sampler.
\textbf{Please check Propositions 2 and 3; the algebra is elementary but I would like
it verified independently before we build on it.}
\end{abstract}

\section{Setting and the point of departure}

Let $\{\Zb_t\}_{t\in\T}$, $\T=[0,T]$, be a L\'evy process with marginal law
$f_u$, and let $\{\xib^{\boldsymbol{x}}_t\}$ be the $\Phib$-initialised L\'evy
random-bridge to $\Psib$ with $\xib^{\boldsymbol{x}}_0=\boldsymbol{x}$. Assume
coordinate-wise independence throughout, so that all statements below are
one-dimensional and apply coordinate-by-coordinate.

The manuscript establishes (Cor.~4.4 and Sec.~4.1) that for
$0\le s<t<T$ the conditional point-to-point bridge
$\xib^{\boldsymbol{x},\boldsymbol{y}}$ has transition density
\begin{equation}
\label{eq:kernel}
\P\big(\xib^{\boldsymbol{x},\boldsymbol{y}}_t\in\d z \,\big|\, \xib^{\boldsymbol{x},\boldsymbol{y}}_s=r\big)
\;=\;\frac{f_{t-s}(z-r)\,f_{T-t}(y-z)}{f_{T-s}(y-r)}\,\d z
\;\triangleq\;\nu_{s,t,T}(r,z,y)\,\d z .
\end{equation}

\begin{rem}[The observation]
\label{rem:main}
The representation involving $\{\Mb_t\}$ is the canonical decomposition of the
bridge into a drift plus a compensated remainder. It is the natural object if one
intends to \emph{integrate an SDE} by a time-stepping scheme, because $\Mb_t$ is
precisely what such a scheme must inject at each step. It is not needed if one
instead \emph{samples \eqref{eq:kernel} directly}. Since \eqref{eq:kernel} is the
exact law of $\xib_t$ given $\xib_s$ and the endpoint, drawing from it advances the
process with no discretisation error whatsoever, for any step size. The martingale
is an artefact of the integrator, not of the process.
\end{rem}

In the generative setting the endpoint $\boldsymbol{y}$ is unknown and is replaced by
the network's estimate $\hat{\boldsymbol{y}}=f^*(\xib_s,s;\Theta^*)\approx
\E[\Yb\mid\xib_s]$. This is the same substitution that DDPM/DDIM ancestral sampling
makes, and it is discussed in Section~\ref{sec:caveat}.

Throughout we use normalised time $u=t/T\in[0,1]$ and write the horizon as $1$.

\section{Gamma random-bridges}

Let $\{Z_t\}$ be a gamma subordinator with $\E[Z_t]=t$ and shape rate $\kappa>0$, so
that increments satisfy $Z_t-Z_s\sim\mathrm{Gamma}\big(\kappa(t-s),\kappa\big)$ with
density
\begin{equation}
f_u(w)=\frac{\kappa^{\kappa u}}{\Gamma(\kappa u)}\,w^{\kappa u-1}e^{-\kappa w},
\qquad w>0 .
\end{equation}

\begin{prop}[Gamma bridge transition]
\label{prop:gamma}
For $0\le s<t<1$ and $r<y$,
\begin{equation}
\xi_t \;=\; \xi_s + (y-\xi_s)\,\beta,
\qquad
\beta\sim\mathrm{Beta}\big(\kappa(t-s),\,\kappa(1-t)\big),
\end{equation}
independently of $\xi_s$. The rate parameter $\kappa$ enters only through the Beta
shape parameters; the scale of the gamma process cancels identically.
\end{prop}

\begin{proof}
Put $a=\kappa(t-s)$, $b=\kappa(1-t)$, so $a+b=\kappa(1-s)$. Substituting the gamma
densities into \eqref{eq:kernel},
\begin{align}
\nu_{s,t,1}(r,z,y)
&=\frac{\dfrac{\kappa^{a}}{\Gamma(a)}(z-r)^{a-1}e^{-\kappa(z-r)}
\cdot \dfrac{\kappa^{b}}{\Gamma(b)}(y-z)^{b-1}e^{-\kappa(y-z)}}
{\dfrac{\kappa^{a+b}}{\Gamma(a+b)}(y-r)^{a+b-1}e^{-\kappa(y-r)}} \notag\\[4pt]
&=\frac{\Gamma(a+b)}{\Gamma(a)\Gamma(b)}\cdot
\frac{(z-r)^{a-1}(y-z)^{b-1}}{(y-r)^{a+b-1}},
\end{align}
the exponential factors cancelling exactly because $(z-r)+(y-z)=(y-r)$, and the
powers of $\kappa$ cancelling because $a+b$ is the total shape. Substituting
$z=r+(y-r)\beta$ with $\d z=(y-r)\d\beta$ and $\beta\in(0,1)$,
\begin{equation}
\nu\,\d z=\frac{\Gamma(a+b)}{\Gamma(a)\Gamma(b)}\,
\beta^{a-1}(1-\beta)^{b-1}\,\d\beta ,
\end{equation}
which is the $\mathrm{Beta}(a,b)$ density.
\end{proof}

\begin{coro}[Training marginal]
With $\xi_0=x$ and terminal $y$, the marginal used for training is
$\xi_s = x+(y-x)\,\beta$, $\beta\sim\mathrm{Beta}(\kappa s,\kappa(1-s))$.
\end{coro}

\begin{coro}[Moments]
\label{coro:gammamoments}
$\E[\xi_t\mid \xi_s,y]=\xi_s+(y-\xi_s)\frac{t-s}{1-s}$ and
\begin{equation}
\mathrm{Var}[\xi_t\mid \xi_s,y]=(y-\xi_s)^2\,
\frac{(t-s)(1-t)}{(1-s)^2\big(\kappa(1-s)+1\big)} .
\end{equation}
\end{coro}

The conditional mean is \emph{identical} to the Gaussian case and reproduces the
linear interpolation property of manuscript eq.~(4.x); only the fluctuation law
differs. As $\kappa\to\infty$ the variance vanishes and the bridge degenerates to
the deterministic linear interpolant.

\begin{rem}[Support]
A gamma subordinator is non-decreasing, so the construction requires
$\xi_s<y$ coordinate-wise. The gamma bridge is therefore a transport between
\emph{ordered} laws, and its paths remain in $[\xi_s,y]$ almost surely. This is the
order-preserving regime motivated in the Introduction of the manuscript.
\end{rem}

\section{Poisson random-bridges}

Let $\{Z_t\}$ be a Poisson process of rate $\lambda$, so
$Z_t-Z_s\sim\mathrm{Poisson}(\lambda(t-s))$ with mass function
$f_u(k)=e^{-\lambda u}(\lambda u)^k/k!$, $k\in\N_0$.

\begin{prop}[Poisson bridge transition]
\label{prop:poisson}
For $0\le s<t<1$ and integers $r\le y$, writing $m=y-r$,
\begin{equation}
\xi_t \;=\; \xi_s + \mathrm{Binomial}\!\left(y-\xi_s,\ \frac{t-s}{1-s}\right).
\end{equation}
The rate $\lambda$ cancels identically.
\end{prop}

\begin{proof}
With $j=z-r$ and $m=y-r$, the discrete analogue of \eqref{eq:kernel} gives
\begin{align}
\P(\xi_t=z\mid \xi_s=r, Y=y)
&=\frac{e^{-\lambda(t-s)}\frac{(\lambda(t-s))^{j}}{j!}\;
        e^{-\lambda(1-t)}\frac{(\lambda(1-t))^{m-j}}{(m-j)!}}
       {e^{-\lambda(1-s)}\frac{(\lambda(1-s))^{m}}{m!}} \notag\\[4pt]
&=\binom{m}{j}\frac{(t-s)^{j}(1-t)^{m-j}}{(1-s)^{m}}
\;=\;\binom{m}{j}p^{j}(1-p)^{m-j},
\end{align}
with $p=\frac{t-s}{1-s}$, since the exponentials combine as
$e^{-\lambda(t-s)}e^{-\lambda(1-t)}=e^{-\lambda(1-s)}$ and the powers of $\lambda$
cancel as $\lambda^{j}\lambda^{m-j}=\lambda^{m}$.
\end{proof}

\begin{coro}
Training marginal: $\xi_s=x+\mathrm{Binomial}(y-x,\,s)$. \\
Moments: $\E[\xi_t\mid\xi_s,y]=\xi_s+(y-\xi_s)\frac{t-s}{1-s}$ and
$\mathrm{Var}[\xi_t\mid\xi_s,y]=(y-\xi_s)\frac{(t-s)(1-t)}{(1-s)^2}$.
\end{coro}

Again the conditional mean is the common linear interpolant. The Poisson bridge is
the classical binomial thinning; the process remains on the integer lattice at every
step, by construction and not by rounding.

\section{The unified sampler}

Writing $\hat y = f^*(\xi_s,s;\Theta^*)$ for the network's endpoint estimate, all
three cases share the update
\begin{equation}
\label{eq:unified}
\xi_u \;=\; \underbrace{\xi_s+(\hat y-\xi_s)\tfrac{u-s}{1-s}}_{\text{common drift}}
\;+\;\underbrace{\text{(driver-specific fluctuation)}},
\end{equation}
with the fluctuation given in Table~\ref{tab:unified}. In every case the fluctuation
degenerates as $u\to1$ and the scheme lands on $\hat y$ exactly, so no boundary
offset $\epsilon$ and no $1/(T-t)$ singularity arise.

\begin{table}[h]
\centering
\begin{tabular}{llll}
\toprule
Driver & Exact transition $\xi_s\to\xi_u$ & Support & Parameter \\
\midrule
Brownian $\sigma W$ &
$\mathcal{N}\!\big(\text{drift},\ \gamma^2\frac{(u-s)(1-u)}{1-s}\big)$ &
$\R$ & $\gamma^2=\sigma^2T$ \\[3pt]
Gamma &
$\xi_s+(\hat y-\xi_s)\,\mathrm{Beta}\big(\kappa(u-s),\kappa(1-u)\big)$ &
$[\xi_s,\hat y]$ & $\kappa$ \\[3pt]
Poisson &
$\xi_s+\mathrm{Bin}\big(\hat y-\xi_s,\ \tfrac{u-s}{1-s}\big)$ &
$\{\xi_s,\dots,\hat y\}\subset\N_0$ & --- \\
\bottomrule
\end{tabular}
\caption{One algorithm, three drivers. The drift is common; only the fluctuation law
changes. Training samples the corresponding marginal from $(x,y)$ in one shot.}
\label{tab:unified}
\end{table}

\begin{rem}[Gaussian case, for completeness]
For $\Zb_t=\sigma\mathbf{W}_t$ the same substitution into \eqref{eq:kernel} gives the
Brownian bridge on $[s,1]$ pinned at $(s,\xi_s)$ and $(1,\hat y)$, i.e.\ the first row
of Table~\ref{tab:unified}. This has been implemented and verified numerically: on a
one-dimensional problem with $\E[Y\mid\xi_s]$ available in closed form, it attains the
Monte-Carlo floor in $1$-Wasserstein distance, whereas the deterministic Euler map
plateaus at $6.9\times$ that floor.
\end{rem}

\section{Caveat: the point-mass substitution}
\label{sec:caveat}

The exactness in Propositions~\ref{prop:gamma} and \ref{prop:poisson} is exactness
\emph{given the endpoint}. In generation we replace the posterior
$\pi_s(\d y)=\P(\Yb\in\d y\mid\xib_s)$ by a point mass at its mean $\hat y$. The
sampler is therefore exact in the SDE-discretisation sense and approximate in the
posterior sense, under-dispersing by $\mathrm{Var}[\Yb\mid\xib_s]$ --- exactly the
quantity Prop.~2.13 of the manuscript shows to be a Loewner-order supermartingale
vanishing at $T$. Two remarks:
\begin{enumerate}
\item This is the same approximation DDPM and DDIM make, so it is not a new
weakness relative to the comparison set.
\item It is removable in principle: a second network head estimating
$\mathrm{Var}[\Yb\mid\xib_s]$ would allow the dropped dispersion to be reinstated.
Prop.~2.13 supplies the theoretical object; whether this helps empirically is open.
\end{enumerate}
A further practical point for the Poisson case: $\hat y$ is real-valued while
$\mathrm{Bin}(\cdot)$ requires an integer count, so $\hat y-\xi_s$ must be rounded or
the head must emit an integer-valued estimate. This is the one place rounding
re-enters and it deserves care.

\section{Proposed head-to-head experiment}

The framework's generality is currently asserted rather than demonstrated, which is
the substance of reviewer point 7. The generality cannot win on MNIST or CIFAR-10:
those are unconstrained real-valued supports, where the Brownian driver is the
correct choice and the framework can at best tie a well-tuned flow-matching or
diffusion model. It can only win where the Gaussian case is \emph{structurally}
wrong.

\paragraph{Claim to target.}
\emph{On count-valued data, a Poisson random-bridge generates on the integer lattice
exactly, at every sampling step, whereas Gaussian diffusion and flow matching must
model a continuous relaxation and round --- and the rounding costs measurably more
at low NFE.}

\paragraph{Design.} A single controlled comparison, all arms sharing architecture,
optimiser, EMA, batch size, training steps and evaluation code; only the transport
changes.
\begin{enumerate}
\item \textbf{Poisson bridge} (ours, Prop.~\ref{prop:poisson}).
\item \textbf{Gaussian bridge} (ours, wrong driver) --- the internal ablation that
isolates the driving process from everything else in the framework.
\item \textbf{Flow matching} with linear path, then rounding --- the modern external
baseline.
\item \textbf{Gaussian diffusion} (DDPM) on the relaxation, then rounding.
\end{enumerate}

\paragraph{Data.} Two tiers.
\emph{(i) Fast and self-contained:} photon-limited images. Take MNIST/CIFAR-10
intensities, fix a mean photon budget (e.g.\ $10$ counts per pixel) and draw Poisson
counts; the task is to generate realistic count images. This is cheap, reproducible
and requires no new data pipeline.
\emph{(ii) Real application:} single-cell RNA-seq count matrices (e.g.\ 10x PBMC).
Counts here are integer, sparse and over-dispersed, and Poisson/negative-binomial is
the \emph{domain-canonical} noise model --- so the Poisson driver is the principled
choice rather than an opportunistic one. Generative modelling of scRNA-seq is an
active area with established metrics.

\paragraph{Metrics.}
\begin{enumerate}
\item \textbf{Support violation rate}: fraction of generated entries outside
$\N_0$ before repair. Ours is identically zero \emph{by construction}; every baseline
is positive. This is structural, not a tuning artefact.
\item \textbf{Distributional fidelity after the baselines' repair step}: FID for
images; for scRNA-seq the standard battery (marginal count histograms, mean--variance
relation, sparsity/zero-inflation rate, nearest-neighbour overlap).
\item \textbf{Quality versus NFE}, since the claimed advantage is in the low-step
regime, reported against wall-clock as well as step count.
\end{enumerate}

\paragraph{Honest risks.} Two, which we should confront rather than discover in
review. First, constraint-respecting continuous methods exist --- Dirichlet and
simplex flow matching, and discrete diffusion (D3PM, SEDD) --- and a credible claim
must include at least one of them, not only vanilla flow matching. Second, if the
rounding penalty turns out to be small, the support-violation argument is
technically correct but uninteresting; we should measure that penalty early and drop
the line of attack if it is negligible.

\section{Numerical verification}

Both propositions were checked against the defining construction --- simulate the
driving process's independent increments, condition on the total, and compare the
conditional law of the first increment with the claimed law. The identities under
test are not used to generate the data. Script: \texttt{notes/verify\_derivations.py},
$4\times10^6$ draws per cell, float64.

\begin{table}[h]
\centering
\small
\begin{tabular}{llrrrrr}
\toprule
& $(\kappa\ \text{or}\ \lambda,\ s,\ t)$ & discrepancy & MC floor & mean & $\frac{t-s}{1-s}$ & var vs theory \\
\midrule
Gamma & $(3,\ 0.5,\ 0.9)$   & KS $0.00063$ & $0.00091$ & $0.79985$ & $0.80000$ & $0.064017 / 0.064000$ \\
Gamma & $(20,\ 0.2,\ 0.5)$  & KS $0.00056$ & $0.00098$ & $0.37499$ & $0.37500$ & $0.013788 / 0.013787$ \\
Gamma & $(20,\ 0.9,\ 0.95)$ & KS $0.00120$ & $0.00109$ & $0.49969$ & $0.50000$ & $0.083338 / 0.083333$ \\
Poisson & $(7,\ 0.5,\ 0.9)$ & TVD $0.00054$ & --- & $3.1995$ & $3.2000$ & $0.6399 / 0.6400$ \\
Poisson & $(30,\ 0.2,\ 0.5)$& TVD $0.00255$ & --- & $8.9997$ & $9.0000$ & $5.6107 / 5.6250$ \\
\bottomrule
\end{tabular}
\caption{Representative cells; all eight gamma and eight Poisson cells tested agree
with theory at the Monte-Carlo floor. The column $\frac{t-s}{1-s}$ confirms the
common drift of \eqref{eq:unified} independently in both families.}
\end{table}

\begin{rem}[Implementation caveat, near-terminal gamma]
When $\kappa(1-t)\ll1$ the Beta law concentrates so sharply at $1$ that
$(\xi_t-\xi_s)/(y-\xi_s)$ rounds to exactly $1$ in floating point --- at $\kappa=3$,
$t=0.999$ this happens in $89.9\%$ of draws. The behaviour is correct (the sampler
lands on the endpoint) but an implementation should special-case the final step
rather than rely on the rounding, as the Gaussian implementation already does.
\end{rem}

\section{Summary of what needs checking}

\begin{enumerate}
\item Proposition~\ref{prop:gamma}: the Beta identity, in particular the claim that
the exponential and $\kappa$ factors cancel identically. (Verified numerically; the
algebra still wants a second pair of eyes.)
\item Proposition~\ref{prop:poisson}: the Binomial identity and the cancellation of
$\lambda$. (Likewise verified numerically.)
\item Remark~\ref{rem:main}: the assertion that sampling $\nu_{s,t,T}$ directly is
legitimate and makes $\{\Mb_t\}$ unnecessary for simulation. This is the load-bearing
claim; everything else is elementary algebra.
\item Whether Corollary~\ref{coro:gammamoments} and its Poisson analogue are
consistent with the conditional-expectation interpolation property already proved in
the manuscript (they appear to be, which is a useful cross-check).
\end{enumerate}

\end{document}
